\documentclass{article}
% Shared layout. Keep benchmark prose and result-specific plots outside this folder.
\PassOptionsToPackage{numbers}{natbib}
\usepackage[preprint]{neurips_2026}
\usepackage[T1]{fontenc}
\usepackage[hidelinks]{hyperref}
\usepackage{booktabs,amsmath,microtype,pgfplots}
\pgfplotsset{compat=1.17}
\usetikzlibrary{patterns}
\definecolor{okblue}{HTML}{0072B2}
\definecolor{okorange}{HTML}{D55E00}
\pgfplotsset{every axis/.append style={font=\small, tick label style={font=\small}, width=\linewidth, height=5cm, grid=major}}

\newcommand{\studyAbstract}{Agreement with a harmful plan and refusal of a reasonable plan are distinct errors. We study calibrated pushback using 12 public plan families and 48 labelled scenarios. Each scenario asks for an ordered decision, a rationale, and an alternative when pushing back. We report local scripted controls and a clearly simulated compliance-bias sweep with global control layers. The label-informed control reaches the scoring ceiling; constant policies do not. A harm-rank shortcut misses labels, while a simulated pushback floor trades lower under-response for higher over-response. These experiments validate score behavior on the released tasks, not deployed model safety. We provide generated data, uncertainty intervals, repeated traces, and a reproducible paper. No learned model or human participant is evaluated.}
\newcommand{\authorOrcid}{0009-0005-9964-6841}
\newcommand{\taskCount}{48}
\newcommand{\planCount}{12}
\newcommand{\conditionCount}{314}
\newcommand{\episodeCount}{30144}
\newcommand{\policyCount}{14}
\newcommand{\adjacentCredit}{0.19}
\newcommand{\courtesyPenalty}{0.25}
\newcommand{\seedCount}{20}
\newcommand{\drawCount}{2000}
\newcommand{\harmMiss}{11}
\newcommand{\oracleScore}{1.000 [1.000, 1.000]}
\newcommand{\oracleExact}{1.000 [0.926, 1.000]}
\newcommand{\harmScore}{0.810 [0.760, 0.865]}
\newcommand{\harmExact}{0.771 [0.635, 0.867]}
\newcommand{\noaltScore}{0.798 [0.713, 0.882]}
\newcommand{\noaltExact}{1.000 [0.926, 1.000]}
\newcommand{\rudeScore}{0.750 [0.750, 0.750]}
\newcommand{\rudeExact}{1.000 [0.926, 1.000]}
\newcommand{\complyScore}{0.298 [0.210, 0.382]}
\newcommand{\complyExact}{0.250 [0.149, 0.388]}
\newcommand{\refuseScore}{0.298 [0.278, 0.317]}
\newcommand{\refuseExact}{0.250 [0.149, 0.388]}
\newcommand{\simnonescore}{0.659 [0.638, 0.681]}
\newcommand{\simnoneunder}{0.365 [0.342, 0.386]}
\newcommand{\simnoneover}{0.000 [0.000, 0.000]}
\newcommand{\simfloorscore}{0.452 [0.439, 0.467]}
\newcommand{\simfloorunder}{0.118 [0.100, 0.134]}
\newcommand{\simfloorover}{0.500 [0.500, 0.500]}
\newcommand{\simrefusescore}{0.298 [0.298, 0.298]}
\newcommand{\simrefuseunder}{0.000 [0.000, 0.000]}
\newcommand{\simrefuseover}{0.750 [0.750, 0.750]}

\title{Calibrated Pushback: A Harm-Ladder Benchmark\\for Sycophantic Agreement}
\author{Pooria Arab\\Independent researcher\\ORCID: \authorOrcid}
\begin{document}
\maketitle
\begin{abstract}\studyAbstract\end{abstract}
\section{Why calibrated pushback matters}
An assistant can agree with a dangerous plan while sounding careful. It can also reject a reasonable plan while sounding responsible. A useful evaluation must distinguish these errors. Agreement alone is not a failure. Refusal alone is not a success. The response must fit the facts and safeguards in the proposed action.
We study this problem through a harm ladder. A plan family keeps a practical goal recognizable while changing its consequences and safeguards. An assistant chooses an ordered response: comply, comply with warning, push back, or refuse. It must also give a rationale. Pushback requires a safer alternative. This structure separates the decision to intervene from the wording used to explain it.
Our narrow construct is \emph{calibrated pushback}: choosing an intervention that matches the labelled risk of the plan as written. We use the benchmark's sycophancy name to describe unjustified agreement with harmful proposals. We do not treat this construct as a complete measure of social sycophancy. The evaluation does not vary a user's identity, authority, prior beliefs, or emotional response. It cannot show that agreement arises from a desire to please.
The contribution is a measurement study of the released harm ladder. We document coverage, inspect the scoring rule through controlled responses, and test a simulated intervention. The controls establish what the score rewards and what simple behavior can earn. The simulated experiment asks whether an indiscriminate pushback floor improves the measured tradeoff. It makes no claim about a trained model's response to a safety prompt.
A central design requirement is to preserve useful agreement. An assistant that blocks every request removes one kind of error by creating another. The benchmark therefore gives proportionate compliance a place in the label space. Its score uses symmetric distance penalties. This symmetry describes calibration against the labels; it does not assert equal real-world costs for every mistake.
\section{Related work and scope}
Sharma et al. study responses that match users' beliefs rather than truth \citep{sharma2023}. Their work connects sycophantic behavior with human preference judgments. Our setting concerns proposed actions and degrees of intervention. It does not estimate how feedback training changes a model. That distinction matters because a poor decision can arise from misunderstanding risk without any social motive.
Perez et al. use model-written evaluations to uncover behaviors, including repetition of a user's preferred answer \citep{perez2022}. The present tasks instead use maintainer-written plans and explicit labels. Their small, inspectable scope supports controlled scoring checks. It offers less breadth than a large behavior survey. Neither authorship method alone establishes that labels reflect expert consensus.
PayBench frames delegated payment behavior around unsafe actions and harmless counterparts \citep{paybench}. Its useful design lesson is that stopping harmful actions and completing permitted actions belong in the same evaluation. Our harm ladders follow that motivation, but they are not matched trap/lookalike pairs. Several facts can change between adjacent variants. We therefore avoid interpreting a score difference as the causal effect of one changed fact.
The companion benchmark-family paper describes this construct and its deterministic grading contract \citep{family}. We retain those definitions. Our contribution is a dedicated analysis of the harm ladder, with directional errors and a simulated control-layer study. We do not import scores, response patterns, or claims from another benchmark in that family.
Evaluation intervals require a stated source of variation \citep{miller2024}. Here, plans form related groups of tasks. We resample complete plans for baseline score intervals. Simulated results use a different unit: random seeds over the fixed public task set. These intervals answer different questions. Neither supplies evidence about human agreement with the labels.
\section{Experiment design}
\subsection{Task and response}
Each task presents a proposed action in ordinary language. The prompt asks whether the user should proceed. The response protocol asks for one structured decision and a concrete rationale. A pushback decision also includes a safer alternative. The experiment runs this protocol through the benchmark's public episode runner and its deterministic scorer.
The benchmark contains \planCount{} plan families and \taskCount{} tasks. Each family has low, medium, high, and severe variants. The labels are not a fixed translation from harm level to response. A medium variant can merit compliance when safeguards make the action proportionate. A high variant can merit a warning when recovery controls reduce its risk. Figure~\ref{fig:coverage} shows this distinction in the released set.
The harm name is descriptive metadata. It is not a substitute for the expected response. In a learned-agent experiment, the assistant should reason from the proposed action and the response protocol. Our harm-rank control deliberately reads that metadata. We label it a shortcut diagnostic, so its result cannot be mistaken for independent language understanding.
\begin{figure}[t]\centering
\begin{tikzpicture}\begin{axis}[ybar stacked,bar width=23pt,ylabel={Task count},xlabel={Stated harm},xtick={0,1,2,3},xticklabels={Low,Medium,High,Severe},legend style={at={(.5,1.23)},anchor=south,legend columns=4},ymin=0]
\addplot[fill=okblue] table[x=harm,y=comply,col sep=comma]{data/coverage.generated.csv};
\addplot[fill=okorange] table[x=harm,y=warn,col sep=comma]{data/coverage.generated.csv};
\addplot[fill=white,pattern=north east lines] table[x=harm,y=pushback,col sep=comma]{data/coverage.generated.csv};
\addplot[fill=gray] table[x=harm,y=refuse,col sep=comma]{data/coverage.generated.csv};
\legend{Comply,Warn,Push back,Refuse}\end{axis}\end{tikzpicture}
\caption{Harm levels do not uniquely determine labels. Counts cover all \taskCount{} tasks in \planCount{} plans. Stacks show expected decisions, not model responses. This is a census of the fixed set; no sampling interval applies.}\label{fig:coverage}\end{figure}

\subsection{Controls and label access}
The baseline battery contains \policyCount{} scripted controls. Empty output, constant text, echo, textual refusal, and an empty first-tool call test failure at the response boundary. The packaged bad policy produces approval text without a valid decision. These controls distinguish informal language from a recorded action. A textual refusal is not equivalent to the valid structured refusal control.
Four other controls always choose one allowed decision. They use the same generic rationale and supply a generic alternative. This removes task interpretation while preserving a valid response. The label-informed packaged control reads the expected decision and label reason. It is an oracle for scoring, not an agent that solves the task. Its ceiling result is necessary for a score check, but it is not evidence of reasoning ability.
The harm-rank shortcut maps successive harm levels onto successive decision bands. Two further oracle controls isolate response requirements. One removes the alternative; the other inserts discourteous wording into the rationale. Both retain the expected decision. These interventions test whether decision accuracy and total score can diverge for the reasons specified by the scorer.
\subsection{Study plan and stopping rule}
We commit the study plan before collecting these results. The plan is a local record, not an independent public registration. Source inspection and existing benchmark checks inform it. That prior knowledge is disclosed because the controls are designed to inspect known scoring rules. The experiment is not a blind prediction of unknown benchmark behavior.
The plan specifies all controls, the simulated grid, exclusions, uncertainty methods, and failure checks. We run every condition twice. Across the study, \conditionCount{} conditions produce \episodeCount{} episodes including repeats. Semantic traces match between repetitions. Repetition establishes determinism for these inputs; it does not create new tasks or increase the inferential sample size.
No observation is excluded because it scores poorly. Invalid responses remain in the denominator and receive zero score. Infrastructure errors stop the run. The experiment uses local scripted adapters and incurs no inference charge. We do not configure a free hosted route, probe provider availability, or evaluate a learned model. Consequently, these results are not a model leaderboard.
\section{Coverage and task interpretation}
\subsection{What a harm ladder changes}
A backup-retention plan illustrates the design. The benign version concerns disposable data that can be rebuilt. Another version concerns internal data with a tested restore path. Riskier versions remove recovery safeguards from production data. The relevant distinction is the combination of value at risk, exposure, reversibility, and available recovery. The word ``backup'' alone cannot determine the response.
A code-review family illustrates why harm ranks and labels can differ. A small reversible change may justify compliance. A more consequential change with tests and rollback can still merit a warning rather than categorical rejection. A proposal that disables safeguards without recovery calls for a stronger response. These are benchmark labels for the stated scenarios, not operational advice for real deployments.
The ladder format also permits within-family analysis. A constant response ignores all changes in the scenario. A rank shortcut responds to the nominal ordering but ignores exceptions created by safeguards. A calibrated response must preserve both the broad trend and those exceptions. This makes the exceptions useful diagnostic items, even when they are a minority of the public set.
\subsection{Coverage is not representativeness}
The public plans span practical decisions, but they do not constitute a probability sample of real requests. The author chooses the scenarios and labels. A census of their distribution is exact for this release. It does not establish prevalence in user conversations, production systems, or professional practice. We therefore report counts as coverage rather than estimates of real-world task frequencies.
The same limitation applies to language and culture. A compact English prompt cannot capture all expectations about authority, politeness, responsibility, or acceptable risk. The task asks for a decision under stated facts. It does not model negotiation, missing context, institutional duties, or a user's response to disagreement. Further validation must test these missing dimensions rather than assuming that more variants solve them.
The present study uses the public set only. A new ordering seed cannot make public plans private. A meaningful hidden evaluation needs new plans and independent labels. Public exposure also permits memorization. A strong score from a future learned model would therefore require held-out evidence before supporting a claim about general decision calibration.
\section{Taxonomy of observed failures}
\paragraph{Under-response.} A valid decision below the expected band gives too little intervention under the label scheme. This includes approval where a warning is expected and warning where refusal is expected. The binary rate does not encode the distance between bands. The total score adds a coarse distance distinction, while raw traces retain the exact decision.
\paragraph{Over-response.} A valid decision above the expected band imposes too much intervention. Blanket refusal is the clearest example. Over-response matters because an assistant should support proportionate plans. Reporting only harmful compliance would reward a system that never helps, even when every proposed action has adequate safeguards.
\paragraph{Missing decision.} Empty text, unstructured approval, or invalid tool arguments do not establish any allowed decision. A missing decision must remain visible. Counting it as safe because no approval was recorded would confuse a protocol failure with successful judgment. We report an invalid rate alongside directional errors.
\paragraph{Unsupported pushback.} The scorer caps pushback that omits an alternative of sufficient length. This is a response-form check. It does not verify that the alternative is feasible or safer. Our generic alternatives make that limit concrete: a valid-length string can satisfy the gate without explaining how to repair the specific plan.
\paragraph{Discourtesy and shortcut success.} The scorer detects a small set of discourteous expressions and reduces credit. This is not a complete measure of respectful communication. Separately, a response can earn substantial credit from harm metadata. Neither a high ordinal correlation nor a valid rationale format proves that the assistant identifies the decisive facts.
\section{Metrics and scoring}
Let $y$ be the expected decision rank and $a$ the recorded decision rank. Ranks follow the order comply, warn, push back, refuse. For a valid decision, define $d=|a-y|$. The base score is $b(d)=1$ for an exact match, $b(d)=\adjacentCredit$ for adjacent bands, and zero otherwise. This rule gives equal base credit to equally distant under-response and over-response.
When pushback lacks an alternative with enough non-whitespace characters, the score is capped at $\adjacentCredit$. When the rationale matches the discourtesy pattern, $\courtesyPenalty$ is subtracted, with a lower bound at zero. These adjustments follow the released scorer. They are design choices rather than fitted estimates of social harm. The study does not change them after observing the controls.
The main result is the arithmetic mean of task scores. Exact-band rate counts a valid recorded rank equal to the expected rank, regardless of the alternative or courtesy adjustment. Under-response and over-response rates count valid lower and higher ranks, respectively. Invalid rate counts missing recorded ranks. All use the full task denominator. Exact, under, over, and invalid therefore partition the outcomes.
This partition prevents a common interpretation error. An oracle can have perfect band accuracy and lose score for its explanation. An empty response can have no under-response and still fail every task. A single scalar hides both distinctions. We report component rates so a reader can determine why the mean changes.
\subsection{Uncertainty and pairing}
Baseline mean-score intervals use \drawCount{} percentile bootstrap draws. Each draw resamples complete plans, keeping their harm variants together. This preserves within-plan dependence. It assumes that these author-selected plans can stand in for an exchangeable plan population, an assumption the study cannot validate. The interval is therefore a sensitivity summary, not a guarantee of coverage for deployment performance.
Binary baseline rates use Wilson intervals under an independent-item approximation. These show boundary uncertainty but do not account for shared plans. We state that limitation rather than presenting narrow component intervals as stronger evidence than the plan-level analysis. Tables retain both kinds of intervals and identify the unit used.
Paired score contrasts subtract the label oracle on the same tasks before resampling plans. Simulated contrasts pair the same seeds and task draws across layers. No significance tests or multiple-comparison claims are made. An interval that collapses at a boundary reflects identical observed values, not certainty about unobserved scenarios. Repeated traces are excluded from every interval's sample size.
\section{Results from scripted controls}
\begin{figure}[t]\centering
\begin{tikzpicture}\begin{axis}[ymin=0,ymax=1.05,ylabel={Mean score},xtick={0,1,2,3,4,5,6,7},xticklabels={Oracle,Comply,Warn,Push back,Refuse,Harm rank,No alternative,Rude},x tick label style={rotate=28,anchor=east}]
\addplot[only marks,mark=*,okblue,error bars/.cd,y dir=both,y explicit] table[x=x,y=score,y error minus=em,y error plus=ep,col sep=comma]{data/baselines.generated.csv};
\end{axis}\end{tikzpicture}\caption{Constant responses remain below the label-informed ceiling. Each point uses \taskCount{} tasks. Bars are percentile intervals from resampling whole plans. Oracle and rude-control intervals collapse because each plan has the same mean. Invalid-output controls appear in Table~\ref{tab:baseline}.}\label{fig:baseline}\end{figure}

\begin{table}[t]\centering\small\begin{tabular}{lrrr}
\toprule
Control & Score & Exact band & Invalid \\
\midrule
Empty & 0.000 [0.000, 0.000] & 0.000 [0.000, 0.074] & 1.000 [0.926, 1.000] \\
Constant text & 0.000 [0.000, 0.000] & 0.000 [0.000, 0.074] & 1.000 [0.926, 1.000] \\
Echo & 0.000 [0.000, 0.000] & 0.000 [0.000, 0.074] & 1.000 [0.926, 1.000] \\
Text refusal & 0.000 [0.000, 0.000] & 0.000 [0.000, 0.074] & 1.000 [0.926, 1.000] \\
First tool & 0.000 [0.000, 0.000] & 0.000 [0.000, 0.074] & 1.000 [0.926, 1.000] \\
Label oracle & 1.000 [1.000, 1.000] & 1.000 [0.926, 1.000] & 0.000 [0.000, 0.074] \\
Packaged bad & 0.000 [0.000, 0.000] & 0.000 [0.000, 0.074] & 1.000 [0.926, 1.000] \\
Always comply & 0.298 [0.210, 0.382] & 0.250 [0.149, 0.388] & 0.000 [0.000, 0.074] \\
Always warn & 0.345 [0.244, 0.446] & 0.250 [0.149, 0.388] & 0.000 [0.000, 0.074] \\
Always push back & 0.345 [0.258, 0.437] & 0.250 [0.149, 0.388] & 0.000 [0.000, 0.074] \\
Always refuse & 0.298 [0.278, 0.317] & 0.250 [0.149, 0.388] & 0.000 [0.000, 0.074] \\
Harm rank & 0.810 [0.760, 0.865] & 0.771 [0.635, 0.867] & 0.000 [0.000, 0.074] \\
Oracle, no alternative & 0.798 [0.713, 0.882] & 1.000 [0.926, 1.000] & 0.000 [0.000, 0.074] \\
Oracle, rude & 0.750 [0.750, 0.750] & 1.000 [0.926, 1.000] & 0.000 [0.000, 0.074] \\
\bottomrule
\end{tabular}
\caption{Scripted controls on \taskCount{} tasks. Score intervals resample plans; rate intervals are Wilson intervals under an independent-item approximation. Brackets show 95\% intervals throughout.}\label{tab:baseline}\end{table}
The label-informed control reaches \oracleScore{}. This is the expected scoring ceiling. Constant comply reaches \complyScore{}, while constant refuse reaches \refuseScore{}. Their equal means conceal opposite errors. The intermediate constant controls also remain below the oracle. The preregistered comparison therefore supports separation from these simple strategies on the released tasks.
The boundary controls all receive no credit. Their failure concerns the structured response protocol. In particular, a refusal sentence is not evidence that the assistant records the correct intervention. A learned model that refuses in free text would fail this interface even if a human considered its wording appropriate. Such a result would require separate analysis of protocol adherence and judgment.
The harm-rank shortcut reaches \harmScore{} and exact-band rate \harmExact{}. It misses \harmMiss{} labels. This supports the preregistered expectation that labels are not a simple harm lookup. At the same time, the shortcut earns substantial credit. That result limits claims that the current public set requires close reading of each proposal.
\begin{figure}[t]\centering
\begin{tikzpicture}\begin{axis}[ymin=0,ymax=3.2,ylabel={Mean decision rank},xtick={0,1,2,3},xticklabels={Low,Medium,High,Severe},legend style={at={(.5,1.02)},anchor=south,legend columns=2}]
\addplot[okblue,mark=*] table[x=harm,y=oracle,col sep=comma]{data/curve.generated.csv};\addplot[okorange,dashed,mark=square*] table[x=harm,y=shortcut,col sep=comma]{data/curve.generated.csv};
\addplot[black,dotted,mark=triangle*] table[x=harm,y=comply,col sep=comma]{data/curve.generated.csv};\addplot[gray,dashdotted,mark=diamond*] table[x=harm,y=refuse,col sep=comma]{data/curve.generated.csv};
\legend{Label oracle,Harm rank,Always comply,Always refuse}\end{axis}\end{tikzpicture}\caption{A rising response curve does not establish label accuracy. Points are fixed-set descriptive means, not estimated population rates. The harm-rank shortcut rises by construction but misses \harmMiss{} labels. Constant extremes expose opposite errors.}\label{fig:curve}\end{figure}

The no-alternative oracle reaches \noaltScore{} despite exact-band rate \noaltExact{}. The rude oracle reaches \rudeScore{} with exact-band rate \rudeExact{}. These controls show that total score measures more than the chosen intervention. They also show why the score should not be described as pure factual accuracy or pure resistance to sycophancy.
The rising harm-rank curve in Figure~\ref{fig:curve} is a further caution. Its ordinal relationship with harm is built into the policy. Constant policies have no decision variation, making their Pearson correlation undefined. A correlation does not replace exact labels, directional errors, or a review of the rationale. We retain the raw correlation in the data but do not use it as a headline measure.
\begin{table}[t]\centering\small\begin{tabular}{lrr}
\toprule
Control & Under-response & Over-response \\
\midrule
Label oracle & 0.000 [0.000, 0.074] & 0.000 [0.000, 0.074] \\
Always comply & 0.750 [0.612, 0.851] & 0.000 [0.000, 0.074] \\
Always warn & 0.500 [0.364, 0.636] & 0.250 [0.149, 0.388] \\
Always push back & 0.250 [0.149, 0.388] & 0.500 [0.364, 0.636] \\
Always refuse & 0.000 [0.000, 0.074] & 0.750 [0.612, 0.851] \\
Harm rank & 0.104 [0.045, 0.222] & 0.125 [0.059, 0.247] \\
Oracle, no alternative & 0.000 [0.000, 0.074] & 0.000 [0.000, 0.074] \\
Oracle, rude & 0.000 [0.000, 0.074] & 0.000 [0.000, 0.074] \\
\bottomrule
\end{tabular}
\caption{Directional error rates for valid-output scripted controls. Brackets show 95\% Wilson intervals with the same item-independence limitation as Table~\ref{tab:baseline}. Invalid controls remain in that table.}\label{tab:directions}\end{table}
\section{Extra virtual experiment: simulated control layers}
\subsection{Simulated population and intervention}
The simulated population starts from the expected decision. For each task and simulated seed, a deterministic pseudorandom draw decides whether to replace that decision with compliance. The replacement probability is $q$. We vary $q$ over the grid shown in Table~\ref{tab:sim}. This creates a controlled spectrum from a label-informed responder to an indiscriminate agreement policy.
This is explicitly a simulation. The starting policy sees the label. The replacement process does not model language generation, social pressure, or learned uncertainty. The simulated parameter is a compliance bias chosen by the experimenter, not a measured psychological trait. We use the construction to inspect score sensitivity and intervention tradeoffs under known assumptions.
We apply three simulated layers to the same underlying draws. The first leaves the decision unchanged. The second raises every lower decision to pushback while retaining stronger decisions. The third replaces every decision with refusal. Each output passes through the actual response environment and scorer. Rationale and alternative text remain generic and valid under the response checks.
The simulated layers intentionally lack semantic context. This isolates the consequence of a global intervention floor. A deployed layer could inspect the task, use uncertainty, or consult a human. Those mechanisms are outside the experiment. A result for this global floor cannot be generalized to every guardrail or every form of safety training.
\subsection{Simulated results and nulls}
\begin{figure}[t]\centering
\begin{tikzpicture}\begin{axis}[ymin=0,ymax=1.05,xlabel={Simulated compliance bias $q$},ylabel={Simulated mean score},legend pos=north east]
\addplot[okblue,mark=*,error bars/.cd,y dir=both,y explicit] table[x=q,y=score,y error minus=em,y error plus=ep,col sep=comma]{data/sim-none.generated.csv};
\addplot[okorange,dashed,mark=square*,error bars/.cd,y dir=both,y explicit] table[x=q,y=score,y error minus=em,y error plus=ep,col sep=comma]{data/sim-floor.generated.csv};
\addplot[black,dotted,mark=triangle*,error bars/.cd,y dir=both,y explicit] table[x=q,y=score,y error minus=em,y error plus=ep,col sep=comma]{data/sim-refuse.generated.csv};
\legend{Simulated: none,Simulated: pushback floor,Simulated: refuse}\end{axis}\end{tikzpicture}\caption{A simulated global floor trades calibration for constraint. Intervals resample \seedCount{} simulated seeds over the fixed tasks. They exclude uncertainty about real agents and new plans. Shared draws pair the simulated control layers.}\label{fig:simulation}\end{figure}

\begin{table}[t]\centering\scriptsize\begin{tabular}{rlrrr}
\toprule
$q$ & Simulated layer & Score & Under-response & Over-response \\
\midrule
0 & none & 1.000 [1.000, 1.000] & 0.000 [0.000, 0.000] & 0.000 [0.000, 0.000] \\
0 & floor & 0.548 [0.548, 0.548] & 0.000 [0.000, 0.000] & 0.500 [0.500, 0.500] \\
0 & refuse & 0.298 [0.298, 0.298] & 0.000 [0.000, 0.000] & 0.750 [0.750, 0.750] \\
0.25 & none & 0.824 [0.806, 0.842] & 0.189 [0.169, 0.208] & 0.000 [0.000, 0.000] \\
0.25 & floor & 0.503 [0.494, 0.510] & 0.055 [0.046, 0.067] & 0.500 [0.500, 0.500] \\
0.25 & refuse & 0.298 [0.298, 0.298] & 0.000 [0.000, 0.000] & 0.750 [0.750, 0.750] \\
0.5 & none & 0.659 [0.638, 0.681] & 0.365 [0.342, 0.386] & 0.000 [0.000, 0.000] \\
0.5 & floor & 0.452 [0.439, 0.467] & 0.118 [0.100, 0.134] & 0.500 [0.500, 0.500] \\
0.5 & refuse & 0.298 [0.298, 0.298] & 0.000 [0.000, 0.000] & 0.750 [0.750, 0.750] \\
0.75 & none & 0.480 [0.461, 0.501] & 0.556 [0.535, 0.577] & 0.000 [0.000, 0.000] \\
0.75 & floor & 0.403 [0.392, 0.416] & 0.178 [0.163, 0.192] & 0.500 [0.500, 0.500] \\
0.75 & refuse & 0.298 [0.298, 0.298] & 0.000 [0.000, 0.000] & 0.750 [0.750, 0.750] \\
1 & none & 0.298 [0.298, 0.298] & 0.750 [0.750, 0.750] & 0.000 [0.000, 0.000] \\
1 & floor & 0.345 [0.345, 0.345] & 0.250 [0.250, 0.250] & 0.500 [0.500, 0.500] \\
1 & refuse & 0.298 [0.298, 0.298] & 0.000 [0.000, 0.000] & 0.750 [0.750, 0.750] \\
\bottomrule
\end{tabular}
\caption{Complete simulated grid. Each cell uses \seedCount{} seeds over the fixed tasks. Brackets show 95\% seed-bootstrap intervals, not uncertainty about a learned-agent population.}\label{tab:sim}\end{table}
At the midpoint of the simulated bias grid, the unchanged simulated policy scores \simnonescore{}. The simulated floor scores \simfloorscore{}. Under-response moves from \simnoneunder{} to \simfloorunder{}, while over-response moves from \simnoneover{} to \simfloorover{}. The floor therefore improves one direction of error at the expense of another under these assumptions.
The simulated refusal layer has score \simrefusescore{} at that midpoint. Its under-response is \simrefuseunder{} and its over-response is \simrefuseover{}. The layer is insensitive to the upstream simulated bias because it overwrites the decision. This invariance is a planned null result. It is not evidence that the simulated agents become equally capable.
The preregistered directional hypothesis needs a boundary qualification. The simulated floor cannot reduce under-response when the unchanged simulated policy already has none. The paired result at that endpoint is null. Elsewhere, the simulated floor removes lower-ranked decisions while retaining some insufficient pushback on refusal-labelled tasks. Appendix~\ref{app:paired} reports the complete paired contrasts, including these boundaries.
The simulated floor's over-response also has no seed variation. Expected compliance and warning labels are always raised too far. A collapsed interval is thus a property of the intervention on this fixed set. It does not mean that the real-world cost of over-response is known exactly. The same distinction applies to the simulated refusal layer's constant results.
\subsection{What the simulated experiment establishes}
The simulated experiment demonstrates a measurement tradeoff, not an improved safety method. A dashboard that showed only under-response would prefer stronger global constraints. The total score and over-response rate reveal lost calibration. This is why both harmful agreement and unnecessary blocking must remain visible when comparing interventions.
A more realistic simulated population could include over-response, context-dependent errors, and missing decisions. Those additions would change the response distribution. They would not turn a simulation into learned-model evidence. The next empirical step is to compare frozen interventions on held-out tasks with real agents and independently reviewed labels. The current result supplies a baseline for that study's design.
\section{Roadmap}
The first priority is independent label review. Reviewers should identify the decisive facts, select a decision band, and record reasons for disagreement. Some scenarios may admit more than one defensible response. An acceptable-label set may be more appropriate than a single rank for those cases. Any change should version the labels and preserve previous results.
The second priority is to separate harm calibration from social pressure. A paired design could hold a plan fixed while changing the user's stated confidence or demand for agreement. Neutral and pressure variants would need equivalent task facts. Only then could a change in response provide stronger evidence of susceptibility to the user's preference rather than difficulty judging the plan.
The third priority is held-out coverage. New plans should vary safeguards independently where feasible. The set should include difficult benign requests as well as obvious hazards. New domains require qualified review before the labels support claims about professional judgment. A larger task count without those checks would increase precision around an uncertain construct.
Finally, future studies should compare context-sensitive controls with the global simulated layers. They should preserve the same task budget, report missing decisions, and show both directional errors. A useful intervention must explain which requests it changes and why. Aggregate improvement without that breakdown would leave its practical effect unclear.
\section{Caveats and responsible interpretation}
This paper reports local scripted experiments and a simulated population sweep. It reports no trained-model evaluation, human baseline, user survey, or field trial. Free hosted inference is not configured for this study, and provider availability is not tested. No statement about a model vendor's relative performance follows from the results.
The oracle uses labels and the rank shortcut uses metadata. Their scores are diagnostics. The scorer accepts generic explanations that pass its length checks. Its courtesy detector is lexical and cannot establish respectful intent. Deterministic replay improves repeatability, but it does not validate the label or the meaning of an explanation. A good replay contract can faithfully reproduce a limited metric.
The response environment records the first valid decision. Later attempts do not necessarily invalidate that state. Our one-turn, one-call design avoids depending on that behavior. These experiments therefore do not validate resistance to repeated-call manipulation or every possible forged trace. Such claims require a separate adversarial study with an explicit failure model.
The fixed public task set shares plan families and may be memorized. The bootstrap assumes exchangeable units, while the Wilson component intervals ignore clustering. Neither addresses label uncertainty. The simulated intervals condition on the task set and the specified error process. We report these limits because a narrow interval can otherwise suggest evidence that was never collected.
Risk also has unequal consequences. A symmetric distance score can assign the same credit to errors with very different real-world outcomes. This paper treats the score as label calibration. It does not estimate expected loss, authorize unsafe actions, or prescribe refusal thresholds for an organization. Deployment decisions need task-specific consequences and accountable human review.
\section{Citation, availability, and authors}
The paper and generated figures use CC-BY-4.0. The benchmark retains its existing MIT license. Code, data, the study plan, and reproduction instructions are available at \url{https://github.com/pooriaarab/benchmark-sycophancy}. The citation metadata identifies this manuscript separately from the software release. This draft has no paper DOI or venue acceptance claim.
Pooria Arab is an independent researcher, ORCID \authorOrcid{}. AI tools assist with implementation and manuscript drafting. The author remains responsible for the labels, claims, and release decision. The study does not claim independent external review. Publication and merge approval remain outside this packet.
\bibliographystyle{plainnat}\bibliography{references}
\appendix
\section{Paired comparisons and interpretation}\label{app:paired}
\begin{table}[h]\centering\small\begin{tabular}{lr}
\toprule
Control & Score minus label oracle \\
\midrule
Empty & -1.000 [-1.000, -1.000] \\
Constant text & -1.000 [-1.000, -1.000] \\
Echo & -1.000 [-1.000, -1.000] \\
Text refusal & -1.000 [-1.000, -1.000] \\
First tool & -1.000 [-1.000, -1.000] \\
Label oracle & 0.000 [0.000, 0.000] \\
Packaged bad & -1.000 [-1.000, -1.000] \\
Always comply & -0.702 [-0.790, -0.618] \\
Always warn & -0.655 [-0.756, -0.554] \\
Always push back & -0.655 [-0.742, -0.563] \\
Always refuse & -0.702 [-0.722, -0.683] \\
Harm rank & -0.190 [-0.240, -0.135] \\
Oracle, no alternative & -0.202 [-0.287, -0.118] \\
Oracle, rude & -0.250 [-0.250, -0.250] \\
\bottomrule
\end{tabular}
\caption{Scripted score differences from the label oracle. Intervals resample whole plans after within-task subtraction. A zero difference for the oracle is an identity check.}\end{table}
The paired baseline comparison preserves the same plans on both sides. It answers whether the control differs from the oracle within the sampled plan groups. Since the oracle supplies expected decisions, this is not a fair competition between agents. It is a check of separation between known response strategies and the score's intended ceiling.
\begin{table}[h]\centering\scriptsize\begin{tabular}{rlrrr}
\toprule
$q$ & Simulated layer & $\Delta$ score & $\Delta$ under & $\Delta$ over \\
\midrule
0 & floor & -0.452 [-0.452, -0.452] & 0.000 [0.000, 0.000] & 0.500 [0.500, 0.500] \\
0 & refuse & -0.703 [-0.703, -0.703] & 0.000 [0.000, 0.000] & 0.750 [0.750, 0.750] \\
0.25 & floor & -0.321 [-0.337, -0.307] & -0.133 [-0.150, -0.116] & 0.500 [0.500, 0.500] \\
0.25 & refuse & -0.527 [-0.545, -0.508] & -0.189 [-0.208, -0.169] & 0.750 [0.750, 0.750] \\
0.5 & floor & -0.207 [-0.222, -0.192] & -0.247 [-0.264, -0.229] & 0.500 [0.500, 0.500] \\
0.5 & refuse & -0.361 [-0.383, -0.341] & -0.365 [-0.386, -0.342] & 0.750 [0.750, 0.750] \\
0.75 & floor & -0.077 [-0.091, -0.064] & -0.378 [-0.393, -0.364] & 0.500 [0.500, 0.500] \\
0.75 & refuse & -0.183 [-0.203, -0.163] & -0.556 [-0.577, -0.535] & 0.750 [0.750, 0.750] \\
1 & floor & 0.048 [0.048, 0.048] & -0.500 [-0.500, -0.500] & 0.500 [0.500, 0.500] \\
1 & refuse & 0.000 [0.000, 0.000] & -0.750 [-0.750, -0.750] & 0.750 [0.750, 0.750] \\
\bottomrule
\end{tabular}
\caption{Simulated paired changes relative to no layer at the same $q$. Brackets show 95\% seed-bootstrap intervals. All simulated boundaries and nulls are retained.}\end{table}
The simulated contrasts retain common random draws. This avoids attributing different simulated mistakes to a control-layer effect. The change in mean score and the changes in directional rates must be read together. A negative under-response contrast can coexist with a negative score contrast because unnecessary intervention also carries a penalty.
\section{Measurement checklist for replication}
A replication should first confirm the task identity and response protocol. It should preserve all harm variants within each plan. If labels change, the study becomes a new version rather than a repeat. If response parsing changes, the investigator should distinguish interface effects from changes in the underlying decisions.
Next, the investigator should replay raw decisions through the scorer and recompute every component rate. A matching scalar mean is insufficient: under-response and over-response can cancel in a summary. The component partition should cover every task exactly once. A missing decision must not disappear when a correlation or curve excludes undefined observations.
For simulated replication, preserve the task draws across layers. Resample seeds only when the target is variation in the simulated population over the fixed set. To study new-plan uncertainty, use a design that samples plans as well. Do not increase the claimed sample size by counting exact repeated runs as independent outcomes.
Finally, compare the prose against the scope of the evidence. A label oracle is not a learned model. A simulated layer is not a deployed safeguard. A length check is not semantic validation. These distinctions are part of the measurement contract and must survive changes to the paper's presentation.
\end{document}
